%% ma: L12-B14-0127
%% lop: 12
%% chuong: 5
%% bai: 14
%% dang: 12.14.1
%% loai: TN4
%% muc: TH
%% dap_an: "B"
%% nguon: "Ảnh giáo viên gửi 25/11/2026 (tài liệu ôn tập Phương trình mặt phẳng), Phần I Câu 5"
%% kien_thuc: "Bài 14 (vectơ pháp tuyến, nhân vectơ với một số khác 0)"
%% trang_thai: chinh-thuc
\begin{ex}
Trong không gian $Oxyz$, cho mặt phẳng $(\alpha)\colon \dfrac{x}{2}-y+\dfrac{z}{5}=1$. Vectơ nào dưới đây là một vectơ pháp tuyến của mặt phẳng $(\alpha)$?
\choice{$\vec n=(-5;1;-2)$}{\True $\vec n=(5;-10;2)$}{$\vec n=(-2;1;-5)$}{$\vec n=(2;-1;5)$}
\loigiai{Một vectơ pháp tuyến của $(\alpha)$ là $\left(\dfrac12;-1;\dfrac15\right)$. Nhân với $10$ được $\vec n=(5;-10;2)$.}
\end{ex}
